\chapter{调用外部 API：RestClient 与 JSON 映射}

这一关写一个 JCDecaux 的客户端：发 HTTP 请求，把返回的 JSON 转成 Java 对象，出错时抛出统一的异常。原 Python 项目里对应的是 \code{main_project/station/services.py} 的 \code{get_all_stations()} 和 \code{format_station()}。

\partinfo{%
  \code{jcdecaux/JcDecauxStationDto.java}：3.1、3.2\\
  \code{jcdecaux/JcDecauxClient.java}：3.3、3.4%
}{\code{./mvnw test -Dtest='Part3*'}}

\section{知识点}

\subsection{JSON 和 record 的转换（Jackson）}

Spring Boot 自带 Jackson，会自动把 JSON 转成 record，前提是\textbf{名字对得上}：

\begin{jsoncode}
{ "title": "Dune", "price": 9.99, "author": { "name": "Herbert" } }
\end{jsoncode}

\begin{java}
public record Book(String title, double price, Author author) {
    public record Author(String name) {}      // 嵌套对象对应嵌套 record
}
\end{java}

JSON 里的字段名和 Java 不一样时（外部 API 经常用 snake\_case），用 \code{@JsonProperty} 指定 JSON 里的名字：

\begin{java}
public record Book(
        String title,
        @JsonProperty("page_count") int pageCount) {}
\end{java}

\begin{note}[包名]
\code{@JsonProperty} 的包名是 \code{com.fasterxml.jackson.annotation}。Spring Boot 4 用的是 Jackson 3，大部分类的包名改成了 \code{tools.jackson}，\textbf{但注解的包名没变}，IntelliJ 自动导入时注意别选错。
\end{note}

\textbf{基本类型和包装类型}：JSON 里某个字段可能是 \code{null} 时，Java 里要用包装类型（\code{Long}、\code{Integer}），因为 \code{long}、\code{int} 不能是 null。看一下测试用的假数据 \code{src/test/resources/fixtures/jcdecaux-stations.json}，第三个站点的 \code{last_update} 就是 \code{null}。

\textbf{时间戳}：JCDecaux 的 \code{last_update} 是毫秒时间戳（从 1970 年开始的毫秒数），用 \code{Instant.ofEpochMilli(long)} 转成 \code{Instant}。

\subsection{DTO 和领域对象}

\code{JcDecauxStationDto} 叫 \textbf{DTO}（Data Transfer Object）：它的结构完全按照 JCDecaux 的 JSON 来，只负责接收数据。我们自己的代码用的是 Part 1 写的 \code{Station}。

为什么要分开？如果 JCDecaux 改了字段名，只需要改 DTO 和转换方法，其他代码都不受影响。

\subsection{RestClient：发 HTTP 请求}

\begin{java}
RestClient client = builder.baseUrl("https://api.example.com").build();

List<Book> books = client.get()                                  // GET 请求
        .uri(uri -> uri.path("/books")                           // 路径拼在 baseUrl 后面
                .queryParam("author", "Herbert")                 // ?author=Herbert
                .build())
        .retrieve()                                              // 发送请求
        .body(new ParameterizedTypeReference<List<Book>>() {});  // 响应体转成 List<Book>
\end{java}

\begin{itemize}
  \item \textbf{为什么注入 \code{RestClient.Builder}，而不是直接 \code{RestClient.create()}？}Spring Boot 提供的 Builder 已经配好了 Jackson 等设置，测试时还能被 \code{MockRestServiceServer} 接管。
  \item \textbf{为什么要 \code{ParameterizedTypeReference}？}Java 的泛型在运行时会被擦除，\code{List<Book>.class} 这种写法不存在。如果只是转成单个对象，直接写 \code{.body(Book.class)} 就行；转成 \code{List<Book>} 这种带泛型的类型，就要用这个写法把完整的类型信息传进去。
\end{itemize}

\subsection{错误处理}

调用外部 API 有两类错误：

\begin{center}
\begin{tabularx}{\textwidth}{X X}
\toprule
情况 & RestClient 的默认行为 \\
\midrule
对方返回 4xx / 5xx & 抛 \code{HttpClientErrorException} / \code{HttpServerErrorException} \\
网络问题：连不上、超时、连接被重置 & 抛 \code{ResourceAccessException} \\
\bottomrule
\end{tabularx}
\end{center}

这些异常都继承自 \code{RestClientException}。我们不希望调用方关心这些细节，所以统一转成自己的 \code{JcDecauxException}（已经写好了）。

\textbf{处理 4xx / 5xx}：在 \code{retrieve()} 之后加 \code{onStatus}：

\begin{java}
.retrieve()
.onStatus(HttpStatusCode::isError, (request, response) -> {
    throw new MyApiException("API returned HTTP " + response.getStatusCode().value());
})
.body(...)
\end{java}

\textbf{处理网络错误}：用 try-catch 把整个调用包起来，捕获 \code{RestClientException}，转成自己的异常。\textbf{一定要把原来的异常作为 cause 传进去}（\code{new MyApiException(message, e)}），否则排查问题时就看不到真正的原因了。

想一想：\code{onStatus} 里抛的 \code{JcDecauxException} 会不会又被外面的 catch 抓住？（不会，因为 \code{JcDecauxException} 不是 \code{RestClientException} 的子类。）

\subsection{测试是怎么验证的}

\code{Part3JcDecauxClientTest} 用了 \code{@RestClientTest(JcDecauxClient.class)}：它只启动和 HTTP 客户端相关的部分，并且把 \code{RestClient.Builder} 接到一个 \code{MockRestServiceServer} 上。测试里写“期望收到某个 URL 的请求，返回某个 JSON”，你的代码发出的请求不会真的到达 JCDecaux。

\section{题目}

\begin{task}{3.1\quad JSON 字段映射}
JCDecaux 返回的一条数据长这样（完整的假数据在 \code{src/test/resources/fixtures/jcdecaux-stations.json}）：

\begin{jsoncode}
{
  "number": 42,
  "contract_name": "dublin",
  "name": "SMITHFIELD NORTH",
  "address": "Smithfield North",
  "position": { "lat": 53.349562, "lng": -6.278198 },
  "banking": false,
  "bonus": false,
  "bike_stands": 30,
  "available_bike_stands": 27,
  "available_bikes": 3,
  "status": "OPEN",
  "last_update": 1761400800000
}
\end{jsoncode}

\code{JcDecauxStationDto} 的字段已经声明好了，但有 5 个字段的 Java 名字和 JSON 名字不一样。给它们加上正确的注解。
\end{task}

\begin{task}{3.2\quad 转换方法}
\textbf{a) \code{lastUpdateInstant()}}：把 \code{lastUpdate}（毫秒时间戳）转成 \code{Instant}；\code{lastUpdate} 为 null 时返回 null。

\begin{example}
\begin{plaincode}
lastUpdate = 1761400800000  ->  2025-10-25T14:00:00Z
lastUpdate = null           ->  null
\end{plaincode}
\end{example}

\textbf{b) \code{toStation()}}：转成领域对象 \code{Station}。注意对应关系：

\medskip
\begin{tabular}{ll}
\toprule
\code{Station} & \code{JcDecauxStationDto} \\
\midrule
\code{availableBikes} & \code{availableBikes} \\
\code{availableStands} & \code{availableBikeStands} \\
\code{totalStands} & \code{bikeStands} \\
\code{lat}、\code{lng} & \code{position.lat}、\code{position.lng} \\
\bottomrule
\end{tabular}
\end{task}

\begin{task}{3.3\quad 发请求}
\textbf{a)} \code{JcDecauxClient} 要交给 Spring 管理，在构造函数里用注入的 \code{RestClient.Builder} 创建 \code{RestClient}，\code{baseUrl} 用配置里的 \code{baseUrl}。

\textbf{b)} 实现 \code{fetchStations()}：

\begin{plaincode}
GET {baseUrl}/stations?contract={contract}&apiKey={apiKey}
\end{plaincode}

查询参数的顺序必须是先 \code{contract} 再 \code{apiKey}（测试会比较完整的 URL）。返回的 JSON 数组转成 \code{List<JcDecauxStationDto>}。

\begin{example}
\begin{plaincode}
baseUrl  = https://jcdecaux.test/vls/v1,  contract = dublin,  apiKey = test-key
请求 URL = https://jcdecaux.test/vls/v1/stations?contract=dublin&apiKey=test-key
\end{plaincode}
\end{example}
\end{task}

\begin{task}{3.4\quad 错误处理}
\begin{itemize}
  \item 对方返回 4xx 或 5xx：抛 \code{JcDecauxException}，消息里要包含状态码（比如 \code{"JCDecaux returned HTTP 403"}）。
  \item 网络错误：抛 \code{JcDecauxException}，并把原始异常作为 cause。
\end{itemize}
\end{task}

\begin{think}
\begin{enumerate}
  \item 外部 API 很慢或者挂了，我们的应用会怎样？应该给 RestClient 设置什么？（提示：连接超时、读取超时）
  \item 如果 JCDecaux 偶尔返回 503，重试是好主意吗？重试几次、间隔多久？什么样的请求不应该重试？
  \item API key 放在 URL 的查询参数里有什么风险？（提示：日志、代理服务器）
\end{enumerate}
\end{think}
